\label{chap:p3-83-25-alpha}

\section[Derivación interna de \texorpdfstring{\(\alpha_{\rm HMT}\)}{alpha HMT}]
{Derivación interna de \texorpdfstring{\(\alpha_{\rm HMT}\)}{alpha HMT}:
estado terminal común, sello dodecafásico y cociclo de acarreo}
\label{p3-83-c25-s1:chap:parte-iii-alpha-hmt}

El generador precedente alcanza un estado terminal enriquecido con dos
publicaciones coordinadas:
\[
\begin{aligned}
w_6\longrightarrow w_{30}\longrightarrow R_{36}\longrightarrow G_9
\longrightarrow x_{\rm term}
&\xrightarrow{\ \pi_{\rm term}\ }B_{\rm term},\\
x_{\rm term}
&\xrightarrow{\ R_{\rm sgn}\ }U_{12}^{\rm sgn}
\xleftrightarrow{\ \mathcal H_{12}\ }K
\longrightarrow A\longrightarrow\alpha_{12}.
\end{aligned}
\]
Las dos flechas terminales publican coordenadas distintas del mismo antecedente.
La transformada \(\mathcal H_{12}\) es reversible; el sistema transversal
\((D_3K,D_4K,Q(K))\) alcanza rango doce; la división euclídea orientada conserva
el cociclo de acarreo y las condiciones de frontera. La cadena utiliza el
estado firmado, el sello y las primitivas declaradas, y obtiene
\(\alpha_{\rm HMT}\) antes de abrir el reconocimiento metrológico.

\subsection{Cierre dodecafásico y construcciones compatibles}
El cierre dodecafásico proporciona el enunciado y la prueba rectora. Las
construcciones E21/22 y \(\mathsf T_{\alpha,\infty}\) mantienen dominios y
funciones propias dentro de la genealogía principal.


\section{Registro firmado, reconstrucción del sello y acarreo integral}
El registro firmado y el cociclo de acarreo despliegan el estado terminal
común, las cartas de Hadamard y diferencias transversales, la forma normal de
Smith, la identidad afín, la prolongación coinductiva y las fronteras
\(662/663\).


\section{Reversibilidad, independencia y no circularidad del cierre dodecafásico}
El estado \(108\), el constructor prospectivo de \(K\), la transformada de
Hadamard bidireccional, el acarreo integral y la prolongación infinita se
articulan mediante controles explícitos de reversibilidad, independencia y no
circularidad. Cada formulación conserva su dominio y su criterio de
refutación.


% Unidad U016. Registro firmado, lectura terminal y reconstruccion de K.
% Redaccion autonoma desde la autoridad material 1063 y su sucesor V2.

\section{Factorización del registro dodecafásico firmado y clausura integral del sello \(K\)}
\label{p3-83-c25-s4:chap:registro-dodecafasico-hadamard-k}

La construcción del sello entero no comienza con \(\alpha\), con su inversa
ni con una tabla metrológica. Su dominio es el estado terminal enriquecido
del Topological Phase Kernel obtenido después de la cadena ya construida
\begin{equation}
 \boxed{
 \mathrm{APP}\longrightarrow\mathrm{TRIT}\longrightarrow\mathrm{TPK}
 \longrightarrow\Omega_{\rm seed}
 \xrightarrow{T_{\min}}\mathcal U_6^{\rm cat}
 \xrightarrow{\rho_3}W_{243}
 \longrightarrow w_{30}\longrightarrow R_{36}
 \longrightarrow G_9\longrightarrow x_{\rm term}.}
 \label{p3-83-c25-s4:eq:cadena-upstream-registro-dodecafasico}
\end{equation}
Las flechas anteriores fueron definidas con sus respectivos dominios:
pares de semillas, firmas, emisiones, palabras ternarias, fronteras,
historias compatibles y holonomía nonádica. Esta unidad no vuelve a escoger
ninguno de esos datos. Construye las cartas enteras del mismo estado
terminal y demuestra su equivalencia.

\subsection{Ventanas de la traza orientada de longitud \(108\)}

Sea
\[
 \Theta_{108}=\mathbb Z/108\mathbb Z
\]
el calendario observable. Sus doce ventanas nonádicas son
\[
 E_m=\{9m+1,\ldots,9m+9\},
 \qquad 0\leq m<12.
\]
La partición es disjunta y exhaustiva:
\[
 \{1,\ldots,108\}=\bigsqcup_{m=0}^{11}E_m.
\]
El conjunto de códigos de dirección que publican un evento se fija como
\[
 D_{\rm evt}=\{2,4,6,8\}.
\]
La orientación de las doce ventanas es
\begin{equation}
 \Sigma_{12}=(+,+,+,-,-,-,+,+,+,-,-,-).
 \label{p3-83-c25-s4:eq:signatura-doce-ventanas}
\end{equation}
Si \(d_t\) es la dirección codificada en el paso \(t\), el evento firmado
de la ventana \(m\) es
\begin{equation}
 e_m=\Sigma_{12}(m)\,
 \#\{t\in E_m:d_t\in D_{\rm evt}\}.
 \label{p3-83-c25-s4:eq:evento-firmado-ventana}
\end{equation}
La fórmula distingue el número no negativo de eventos y el signo de la
orientación. No reemplaza una dirección cardinal por un signo: las dos
coordenadas siguen presentes en el estado TPK.

\begin{definition}[Registro enriquecido dodecafásico]
\label{p3-83-c25-s4:def:registro-enriquecido-dodecafasico}
El registro transportado por la traza es
\[
 \mathcal R_{12}(x)=
 \bigl((E_m,\tau_m,\sigma_m,\kappa_m,\Lambda_m)\bigr)_{m=0}^{11},
\]
donde \(\tau_m\) es la orientación TRIT, \(\sigma_m\) la firma de
frontera, \(\kappa_m\) el acarreo, y \(\Lambda_m\) el segmento de registro genealógico
que conserva la ruta. Existe una cadena de proyecciones
\[
 \mathcal X_{\rm TPK}^{\rm enr}
 \longrightarrow\mathcal R_{12}
 \longrightarrow\Theta_{108},
\]
pero ninguna de las dos flechas es una identificación.
\end{definition}

En particular, \(\Theta_{108}\) sólo conserva la posición del calendario.
El registro \(\mathcal R_{12}\) conserva además orientación, firma,
acarreo y procedencia. El estado completo conserva todavía cilindro,
frontera, cociente, terminal y memoria vertical.

\subsection{Carta integral de la oposición dodecafásica}

Para \(0\leq i<6\), las posiciones opuestas determinan
\begin{equation}
 p_i=e_i+e_{i+6},
 \qquad
 a_i=e_i-e_{i+6}.
 \label{p3-83-c25-s4:eq:pares-opuestos-registro}
\end{equation}
La carga central y los cinco márgenes relativos son
\begin{equation}
 s_C=\sum_{i=0}^{5}p_i,
 \qquad
 M_i=p_i-p_5\quad(0\leq i<5).
 \label{p3-83-c25-s4:eq:carga-margenes-registro}
\end{equation}

\begin{definition}[Carta integral \(1+5+6\)]
\label{p3-83-c25-s4:def:carta-integral-156}
Se define
\begin{equation}
 \mathcal L_{12}(e_0,\ldots,e_{11})
 =\bigl(s_C,M_0,\ldots,M_4,a_0,\ldots,a_5\bigr).
 \label{p3-83-c25-s4:eq:carta-integral-156}
\end{equation}
La partición de coordenadas es
\[
 1+5+6=12.
\]
\end{definition}

\begin{theorem}[Inversión exacta de la carta integral]
\label{p3-83-c25-s4:thm:inversion-carta-integral-156}
Una tupla
\[
 \ell=(s_C,M_0,\ldots,M_4,a_0,\ldots,a_5)\in\mathbb Z^{12}
\]
pertenece a la imagen de \(\mathcal L_{12}\) si y sólo si
\begin{align}
 s_C-\sum_{i=0}^{4}M_i&\equiv0\pmod6,
 \label{p3-83-c25-s4:eq:congruencia-seis-carta}\\
 p_i&\equiv a_i\pmod2
 \quad(0\leq i<6),
 \label{p3-83-c25-s4:eq:congruencia-paridad-carta}
\end{align}
donde
\begin{align}
 p_5&=\frac{s_C-\sum_{i=0}^{4}M_i}{6},
 \label{p3-83-c25-s4:eq:inversa-p5-carta}\\
 p_i&=p_5+M_i\quad(0\leq i<5).
 \label{p3-83-c25-s4:eq:inversa-pi-carta}
\end{align}
En ese caso, la única preimagen entera es
\begin{equation}
 e_i=\frac{p_i+a_i}{2},
 \qquad
 e_{i+6}=\frac{p_i-a_i}{2}.
 \label{p3-83-c25-s4:eq:inversa-eventos-carta}
\end{equation}
\end{theorem}

\begin{proof}
Sumando \(M_i=p_i-p_5\) para \(0\leq i<5\), obtenemos
\[
 \sum_{i=0}^{4}M_i
 =\sum_{i=0}^{4}p_i-5p_5
 =s_C-6p_5,
\]
lo que fuerza \eqref{p3-83-c25-s4:eq:inversa-p5-carta} y la congruencia módulo seis.
Una vez reconstruidos los seis \(p_i\), el sistema
\[
 p_i=e_i+e_{i+6},
 \qquad a_i=e_i-e_{i+6}
\]
tiene la única solución \eqref{p3-83-c25-s4:eq:inversa-eventos-carta}. Esa solución es
entera exactamente cuando \(p_i\) y \(a_i\) tienen la misma paridad.
Las dos familias de congruencias son, por tanto, necesarias y suficientes.
\end{proof}

Este teorema es anterior al cierre numérico de \(\alpha\). Explica cómo la
traza orientada produce una carta integral reversible, y por qué no basta
con conservar el calendario reducido.

\subsection{Selección del estado terminal común}

Las tres secciones compatibles llegan a profundidad terminal con catálogos
de cardinalidades
\[
 |\mathcal T_\pi|=196,
 \qquad
 |\mathcal T_e|=197,
 \qquad
 |\mathcal T_\varphi|=196.
\]
El producto contiene
\begin{equation}
 196\cdot197\cdot196=7\,567\,952
 \label{p3-83-c25-s4:eq:cardinal-producto-terminal}
\end{equation}
ternas ordenadas.

El margen de columnas del lector terminal es
\[
 q_\partial=(1,2,2,1,2,0)\in\mathbb F_3^6,
\]
y la orientación de filas del canal conjunto es
\[
 \sigma=(1,1,-1)=(1,1,2)\in\mathbb F_3^3.
\]
El primer dato actúa sobre columnas; el segundo sobre filas. No se deduce
uno del otro.

Las firmas locales de Gram, norma, peso y distancia reducen
\eqref{p3-83-c25-s4:eq:cardinal-producto-terminal} a \(331\) ternas. El margen
\(q_\partial\) deja dos índices:
\[
 (120,197,157),
 \qquad
 (129,197,148).
\]
Sus matrices son
\[
 B_0=
 \begin{pmatrix}
 0&2&0&2&2&0\\
 0&0&1&2&2&0\\
 1&0&1&0&1&0
 \end{pmatrix},
 \qquad
 B_1=
 \begin{pmatrix}
 0&2&1&0&2&0\\
 0&0&1&2&2&0\\
 1&0&0&2&1&0
 \end{pmatrix}.
\]
Las sumas de filas son
\[
 (0,2,0),
 \qquad
 (2,2,1).
\]
Como
\[
 \langle\sigma\rangle
 =\{(0,0,0),(1,1,2),(2,2,1)\},
\]
sólo \(B_1\) satisface la condición axial.

\begin{theorem}[Unicidad de la publicación terminal visible]
\label{p3-83-c25-s4:thm:unicidad-publicacion-terminal-visible}
La selección conjunta determina
\begin{equation}
 B_{\rm term}=
 \begin{pmatrix}
 021020\\001220\\100210
 \end{pmatrix},
 \label{p3-83-c25-s4:eq:matriz-terminal-visible}
\end{equation}
correspondiente a la terna \((129,197,148)\). Ninguna evaluación decimal de
\(\alpha\) interviene en el censo ni en las dos condiciones de selección.
\end{theorem}

\begin{proof}
El censo y el margen dejan exactamente \(B_0,B_1\). La carga de \(B_0\)
no pertenece a \(\langle\sigma\rangle\), mientras la de \(B_1\) es
\(2\sigma\). La segunda candidata es, por tanto, la única que satisface
simultáneamente ambos requisitos tipados.
\end{proof}

\subsection{2 publicaciones del mismo estado y ausencia de conversión retrospectiva}

Sea
\[
 x_{\rm term}\in
 \mathcal X_{\rm TPK}^{\rm term,enr}
 \subseteq\mathcal X_{\rm TPK}^{[108]}.
\]
La lectura visible es
\[
 \pi_{\rm term}(x_{\rm term})=B_{\rm term}.
\]
La carta firmada conserva otra coordenada del mismo estado.

\begin{definition}[Lectura causal firmada terminal]
\label{p3-83-c25-s4:def:subespacio-firmado-terminal}
El dominio no incorpora como coordenadas de entrada ni \(K\) ni \(U\). El
estado \(x_{\rm term}\), producido por la composición completa
\(\mathrm{APP}\to\mathrm{TRIT}\to\mathrm{TPK}\), contiene el registro
terminal enriquecido, y el extractor de sus doce ventanas orientadas es
\begin{equation}
 R_{12}:\mathcal X_{\rm TPK}^{\rm term,enr}
 \longrightarrow\mathcal R_{12}^{\rm or},
 \qquad
 R_{12}(x_{\rm term})
 =\bigl((E_m,\tau_m,\sigma_m,\kappa_m,\Lambda_m)\bigr)_{m=0}^{11}.
 \label{p3-83-c25-s4:eq:subespacio-firmado-terminal}
\end{equation}
El lector de canales \(90/120\) del ledger emite, sin recibir el resultado
como dato,
\[
 E_{108}^{90,120}:\mathcal R_{12}^{\rm or}
 \longrightarrow\mathbb Z^{12}\times\mathbb Z^{12}\times\mathbb Z,
 \qquad
 E_{108}^{90,120}\bigl(R_{12}(x_{\rm term})\bigr)
 =(b_{90},b_{120},Q_{\rm TPK}),
\]
con
\begin{align*}
 b_{90}={}&(-495,-116,-684,108,601,-90,
             -173,-88,313,560,-397,461),\\
 b_{120}={}&(-425,-281,-481,671,-255,30,
              475,-543,680,251,6,-128),\\
 Q_{\rm TPK}={}&6263.
\end{align*}
Para \(r=1,2,3\) y \(j=0,1,2,3\), se define
\[
\begin{aligned}
 B_r&=\sum_{j=0}^{3}b_{120,r+3j},
 &d_{r,j}&=b_{90,r+3j},\\
 A_1&=\frac{Q_{\rm TPK}+2B_1+B_2}{3},
 &A_2&=A_1-B_1,
 &A_3&=A_2-B_2.
\end{aligned}
\]
La publicación \(\Pi_H\) forma los tres bloques
\[
 U^{(r)}=
 \bigl(A_r,d_{r,1}-d_{r,3},d_{r,0}+d_{r,2},
 d_{r,0}-d_{r,2}\bigr),
\]
y los reintercala en el orden dodecafásico, es decir,
\[
 \Pi_H(b_{90},b_{120},Q_{\rm TPK})
 :=\operatorname{reint}_3\bigl(U^{(1)},U^{(2)},U^{(3)}\bigr)=U.
\]
Sea \(\mathcal H_{12}\) la
transformación integral que se construirá en la sección siguiente. El módulo
de publicación es
\begin{equation}
 \mathcal U_{12}^{\rm sgn}
 :=\operatorname{im}
 \bigl(\mathcal H_{12}:\mathbb Z^{12}\to\mathbb Z^{12}\bigr).
 \label{p3-83-c25-s4:eq:modulo-publicacion-firmada}
\end{equation}
La lectura firmada es, causalmente, la composición
\begin{equation}
 R_{\rm sgn}:=
 \Pi_H\circ E_{108}^{90,120}\circ R_{12}:
 \mathcal X_{\rm TPK}^{\rm term,enr}
 \longrightarrow\mathcal U_{12}^{\rm sgn},
 \qquad
 R_{\rm sgn}(x_{\rm term})=U.
 \label{p3-83-c25-s4:eq:lector-firmado-terminal}
\end{equation}
Sólo después de esta salida se reconstruye
\(K=\mathcal H_{12}^{-1}U\); la identidad
\(U=\mathcal H_{12}K\) certifica esa reconstrucción y no define el dominio
del lector.
\end{definition}

La escritura conjunta correcta es
\begin{equation}
 x_{\rm term}
 \xrightarrow{(\pi_{\rm term},R_{\rm sgn})}
 (B_{\rm term},U_{12}^{\rm sgn}).
 \label{p3-83-c25-s4:eq:doble-publicacion-terminal}
\end{equation}
No se interpone una flecha
\(B_{\rm term}\to U_{12}^{\rm sgn}\): la matriz visible ha olvidado
precisamente las coordenadas de orientación, oposición, cociente, acarreo y
registro genealógico utilizadas por la carta firmada. La fuente de ambas publicaciones es
el estado enriquecido producido por \eqref{p3-83-c25-s4:eq:cadena-upstream-registro-dodecafasico}.

La evaluación del registro firmado da
\begin{equation}
 \boxed{
 U=(2378,1406,2479,-452,998,-551,
 -668,-204,-371,-322,-28,-997).}
 \label{p3-83-c25-s4:eq:vector-u-firmado-autonomo}
\end{equation}

\subsection{Naturalidad respecto de la profundidad}

Para \(N'\geq N\geq108\), los truncamientos del sistema firmado son
\begin{equation}
 \tau_{N',N}^{\rm sgn}
 x_{\rm term}^{[N']}
 =\tau_{N',N}x_{\rm term}^{[N']}.
 \label{p3-83-c25-s4:eq:truncamiento-preserva-ku-autonomo}
\end{equation}
Por tanto, el truncamiento actúa sólo sobre la historia y el estado
enriquecido; no transporta como entradas independientes ni \(K\) ni \(U\).

\begin{theorem}[Cuadrado de naturalidad de la lectura firmada]
\label{p3-83-c25-s4:thm:naturalidad-lector-firmado-autonomo}
Se cumple
\begin{equation}
 \boxed{
 R_{{\rm sgn},N}\circ\tau_{N',N}^{\rm sgn}
 =\operatorname{id}_{\mathcal U_{12}^{\rm sgn}}
  \circ R_{{\rm sgn},N'}.}
 \label{p3-83-c25-s4:eq:cuadrado-naturalidad-lector-firmado}
\end{equation}
\end{theorem}

\begin{proof}
La naturalidad de los estados conserva literalmente las primeras ciento ocho
ventanas del ledger:
\[
 R_{12,N}\circ\tau_{N',N}^{\rm sgn}=R_{12,N'}.
\]
El truncamiento puede modificar sólo la historia ulterior. Como
\(E_{108}^{90,120}\) y \(\Pi_H\) dependen exclusivamente de esas ventanas,
ambos miembros devuelven el mismo \(U\). La conmutatividad es una igualdad de
aplicaciones, no una analogía entre diagramas.
\end{proof}

\subsection{Órbitas de paso \(3\) y transformada de Hadamard}

El orden cíclico del sello se distribuye en las tres órbitas
\begin{align}
 K^{(1)}&=(K_1,K_4,K_7,K_{10}),
 \label{p3-83-c25-s4:eq:orbita-k-1}\\
 K^{(2)}&=(K_2,K_5,K_8,K_{11}),
 \label{p3-83-c25-s4:eq:orbita-k-2}\\
 K^{(3)}&=(K_3,K_6,K_9,K_{12}).
 \label{p3-83-c25-s4:eq:orbita-k-3}
\end{align}
Sea
\begin{equation}
 H_4=
 \begin{pmatrix}
 1&1&1&1\\
 1&1&-1&-1\\
 1&-1&1&-1\\
 1&-1&-1&1
 \end{pmatrix}.
 \label{p3-83-c25-s4:eq:matriz-hadamard-cuatro-autonoma}
\end{equation}
Si \(P_{\rm orb}\) reordena las doce coordenadas de acuerdo con
\eqref{p3-83-c25-s4:eq:orbita-k-1}--\eqref{p3-83-c25-s4:eq:orbita-k-3}, la transformación global es
\begin{equation}
 \mathcal H_{12}
 =P_{\rm orb}^{-1}
 (H_4\oplus H_4\oplus H_4)P_{\rm orb}.
 \label{p3-83-c25-s4:eq:hadamard-global-doce}
\end{equation}

Los tres bloques de \eqref{p3-83-c25-s4:eq:vector-u-firmado-autonomo}, leídos por
órbitas, son
\begin{align*}
 U^{(1)}&=(2378,-452,-668,-322),\\
 U^{(2)}&=(1406,998,-204,-28),\\
 U^{(3)}&=(2479,-551,-371,-997).
\end{align*}

\begin{lemma}[Cuarto de vuelta algebraico de Hadamard]
\label{p3-83-c25-s4:lem:cuarto-inversa-hadamard}
La matriz \(H_4\) satisface
\begin{equation}
 H_4^2=4I_4,
 \qquad
 H_4^{-1}=\frac14H_4,
 \qquad
 \det H_4=-16.
 \label{p3-83-c25-s4:eq:identidades-hadamard-cuatro}
\end{equation}
\end{lemma}

\begin{proof}
Las filas de \(H_4\) son ortogonales y cada una posee norma cuadrada cuatro.
Como la matriz es simétrica, \(H_4^{\mathsf T}H_4=H_4^2=4I_4\). La
fórmula de la inversa y el determinante se deducen inmediatamente.
\end{proof}

\begin{theorem}[Reconstrucción entera única del sello]
\label{p3-83-c25-s4:thm:reconstruccion-entera-k-hadamard}
La inversión
\begin{equation}
 K^{(r)}=\frac14H_4U^{(r)}
 \qquad(r=1,2,3)
 \label{p3-83-c25-s4:eq:inversion-bloques-hadamard}
\end{equation}
produce
\begin{align*}
 K^{(1)}&=(234,729,621,794),\\
 K^{(2)}&=(543,659,058,146),\\
 K^{(3)}&=(140,824,914,601).
\end{align*}
Al deshacer el orden orbital se obtiene
\begin{equation}
 \boxed{
 K=(234,543,140,729,659,824,621,058,914,794,146,601).}
 \label{p3-83-c25-s4:eq:sello-k-doce-autonomo}
\end{equation}
Es la única preimagen racional de \(U\), y pertenece a
\(\mathbb Z^{12}\).
\end{theorem}

\begin{proof}
Por \eqref{p3-83-c25-s4:eq:identidades-hadamard-cuatro}, la preimagen racional es única.
Las multiplicaciones explícitas son
\begin{align*}
 H_4U^{(1)}&=(936,2916,2484,3176),\\
 H_4U^{(2)}&=(2172,2636,232,584),\\
 H_4U^{(3)}&=(560,3296,3656,2404).
\end{align*}
Todas las coordenadas son divisibles por cuatro. Dividir y reintercalar
produce \eqref{p3-83-c25-s4:eq:sello-k-doce-autonomo}.
\end{proof}

El factor \(1/4\) no se escoge para ajustar una salida: es la inversa
forzada por \(H_4^2=4I_4\). Tampoco es la raíz cuadrada de \(-1\). La
estructura compleja interna \(A_W^2=-I\) aparecerá después de construir la
matriz de Paley; son dos hechos de tipos diferentes.

\subsection{Estructura integral y forma normal de Smith}

\begin{theorem}[Imagen integral de \(H_4\)]
\label{p3-83-c25-s4:thm:smith-hadamard-autonomo}
La forma normal de Smith es
\begin{equation}
 \operatorname{SNF}(H_4)=\operatorname{diag}(1,2,2,4).
 \label{p3-83-c25-s4:eq:snf-hadamard-autonoma}
\end{equation}
En consecuencia,
\begin{equation}
 \operatorname{coker}\mathcal H_{12}
 \cong(\mathbb Z/2\mathbb Z)^6
 \oplus(\mathbb Z/4\mathbb Z)^3,
 \label{p3-83-c25-s4:eq:cociente-hadamard-doce}
\end{equation}
y la imagen de \(\mathcal H_{12}\) posee índice
\[
 16^3=4096
\]
en \(\mathbb Z^{12}\). Para \(u\in\mathbb Z^4\),
\begin{equation}
 u\in H_4\mathbb Z^4
 \quad\Longleftrightarrow\quad
 H_4u\equiv0\pmod4.
 \label{p3-83-c25-s4:eq:criterio-imagen-hadamard}
\end{equation}
\end{theorem}

\begin{proof}
El máximo común divisor de las entradas es uno. Los menores de orden dos son
pares y alguno tiene valor absoluto dos; los de orden tres son múltiplos de
cuatro y alguno tiene valor absoluto cuatro; el determinante tiene módulo
dieciséis. Los divisores determinantes son \(1,2,4,16\), de donde se obtienen
los factores invariantes \(1,2,2,4\). La suma directa triple repite los
factores y multiplica los índices. Finalmente,
\(H_4^{-1}=H_4/4\), por lo que la preimagen es entera exactamente bajo la
congruencia \eqref{p3-83-c25-s4:eq:criterio-imagen-hadamard}.
\end{proof}

\subsection{Certificación y reconstrucción por diferencias de pasos \(3\) y \(4\)}

Una vez que el ledger ha emitido \((b_{90},b_{120},Q_{\rm TPK})\), que
\(\Pi_H\) ha producido \(U\) y que Hadamard ha reconstruido \(K\), las
siguientes coordenadas actúan aguas abajo como certificado independiente de
esas salidas previas. Con índices cíclicos módulo doce, definimos
\begin{equation}
 (D_3K)_m=K_m-K_{m+3},
 \qquad
 (D_4K)_m=K_m-K_{m+4},
 \qquad
 Q(K)=\sum_{m=1}^{12}K_m.
 \label{p3-83-c25-s4:eq:extractor-diferencias-k}
\end{equation}
Para \eqref{p3-83-c25-s4:eq:sello-k-doce-autonomo}, se verifica
\begin{align*}
 D_3K=b_{90}={}&(-495,-116,-684,108,601,-90,\\
         &\qquad-173,-88,313,560,-397,461),\\
 D_4K=b_{120}={}&(-425,-281,-481,671,-255,30,\\
         &\qquad475,-543,680,251,6,-128),\\
 Q(K)=Q_{\rm TPK}={}&6263.
\end{align*}

\begin{theorem}[Completitud del extractor transversal]
\label{p3-83-c25-s4:thm:completitud-extractor-transversal-k}
El operador
\[
 \mathcal E_{\rm dod}=(D_3,D_4,Q):
 \mathbb Z^{12}\longrightarrow\mathbb Z^{25}
\]
tiene rango doce y forma normal de Smith
\begin{equation}
 \operatorname{SNF}(\mathcal E_{\rm dod})
 =\operatorname{diag}(\underbrace{1,\ldots,1}_{11},12).
 \label{p3-83-c25-s4:eq:snf-extractor-transversal-k}
\end{equation}
En particular, \((D_3K,D_4K,Q(K))\) determina un único \(K\). Esta
unicidad certifica y reconstruye los datos ya emitidos por el ledger; no los
genera ni los selecciona.
\end{theorem}

\begin{proof}
Si \(D_3v=D_4v=0\), entonces \(v\) es constante a lo largo de los pasos
tres y cuatro. Como \(\gcd(3,4)=1\), esos pasos generan todo
\(\mathbb Z/12\mathbb Z\), de modo que
\[
 \ker D_3\cap\ker D_4=\langle\mathbf1\rangle.
\]
La carga total no se anula en esa recta, pues \(Q(\mathbf1)=12\). El rango
es doce. Las filas de \((D_3,D_4)\) son incidencias orientadas de un grafo
de Cayley conexo; un árbol generador aporta once factores invariantes
unitarios. Todo menor máximo debe contener la fila \(Q\), y las operaciones
unimodulares reducen su última contribución a doce. Existe un menor máximo de
módulo doce; por tanto el último factor es exactamente doce.
\end{proof}

La reconstrucción certificadora de \(U\) desde estas coordenadas reproduce la
publicación previa de \(\Pi_H\) y no necesita volver a multiplicar por
\(H_4\). Para \(r=1,2,3\), sea
\begin{equation}
 B_r=\sum_{j=0}^{3}(D_4K)_{r+3j}
 =\sum_{j=0}^{3}b_{120,r+3j}.
 \label{p3-83-c25-s4:eq:br-diferencias-k}
\end{equation}
Entonces
\[
 (B_1,B_2,B_3)=(972,-1073,101).
\]
Las sumas orbitales son
\begin{align}
 A_1&=\frac{Q+2B_1+B_2}{3}=2378,
 \label{p3-83-c25-s4:eq:a1-reconstruccion-u}\\
 A_2&=A_1-B_1=1406,
 \label{p3-83-c25-s4:eq:a2-reconstruccion-u}\\
 A_3&=A_2-B_2=2479.
 \label{p3-83-c25-s4:eq:a3-reconstruccion-u}
\end{align}
Si \(d_{r,j}=(D_3K)_{r+3j}=b_{90,r+3j}\), el bloque firmado de la órbita
\(r\) es
\begin{equation}
 \bigl(A_r,
 d_{r,1}-d_{r,3},
 d_{r,0}+d_{r,2},
 d_{r,0}-d_{r,2}\bigr),
 \label{p3-83-c25-s4:eq:bloque-firmado-desde-diferencias}
\end{equation}
que reproduce los tres bloques \(U^{(r)}\) ya publicados por \(\Pi_H\).

\subsection{Tétrada de umbral del sello}

La completación decimal local satisface
\[
 1000=729+271.
\]
El soporte de las coordenadas del sello que alcanzan el umbral \(729\) es
\begin{equation}
 \boxed{
 B_K:=\{m:K_m\geq729\}=\{4,6,9,10\}.}
 \label{p3-83-c25-s4:eq:tetrada-umbral-k}
\end{equation}
Esta tétrada se obtiene después de construir \(K\). En esta unidad no se la
denomina todavía estrella de Witt: tal objeto sólo estará definido después
de construir el código ternario, sus hexadas y el sistema de Steiner.

\subsection{Independencia causal respecto de la salida física}

\begin{theorem}[Aislamiento de \(U\) y \(K\)]
\label{p3-83-c25-s4:thm:aislamiento-u-k-alpha-codata}
La composición
\begin{equation}
 \Omega_{\rm seed}\longrightarrow x_{\rm term}
 \xrightarrow{R_{\rm sgn}}U
 \xrightarrow{\mathcal H_{12}^{-1}}K
 \label{p3-83-c25-s4:eq:composicion-upstream-u-k}
\end{equation}
no recibe como entrada \(\alpha\), \(\alpha^{-1}\), CODATA ni una cadena
decimal objetivo. El sello \(K\) queda determinado antes de definir la
recurrencia que producirá \(\alpha\).
\end{theorem}

\begin{proof}
La primera parte de la composición utiliza únicamente semillas, emisiones,
lectores ternarios, operadores \(L_0,L_1\), relaciones de supervivencia,
transporte de fase y las cartas del estado TPK. La lectura
\(R_{\rm sgn}=\Pi_H\circ E_{108}^{90,120}\circ R_{12}\) es la composición
causal del registro terminal, y la última flecha es la transformación lineal
\(\mathcal H_{12}^{-1}=\mathcal H_{12}/4\)
sobre su imagen integral. Ninguna de las fórmulas contiene una variable
física o una tabla metrológica. La recurrencia en base mil se introduce sólo
en la unidad siguiente, después de haber fijado \(K\).
\end{proof}

\subsection{Obstrucciones del cierre dodecafásico: sello, acarreo, reversibilidad y no circularidad}

\begin{remark}[Certificado y control de validez: Refutación del registro y del sello]
La construcción queda refutada si se verifica cualquiera de los hechos
siguientes:
\begin{enumerate}
 \item las ventanas \(E_m\) no particionan los ciento ocho pasos;
 \item la carta \(\mathcal L_{12}\) incumple su inversa o sus congruencias;
 \item el censo terminal no produce \(331\) ternas, dos matrices y una única
       carga axial;
 \item se intenta obtener \(U\) de \(B_{\rm term}\) después de olvidar las
       coordenadas firmadas;
 \item falla el cuadrado de naturalidad
       \eqref{p3-83-c25-s4:eq:cuadrado-naturalidad-lector-firmado};
 \item \(H_4^2\neq4I_4\), alguna preimagen no es entera o el vector
       reconstruido difiere de \eqref{p3-83-c25-s4:eq:sello-k-doce-autonomo};
 \item la forma normal de Smith difiere de
       \eqref{p3-83-c25-s4:eq:snf-hadamard-autonoma};
 \item el certificado \((D_3,D_4,Q)\) pierde rango, no coincide con las
       emisiones previas del ledger o no reproduce \(U\);
 \item \(B_K\) se introduce antes de construir \(K\), o se le atribuye una
       estructura de Witt antes de definir \(W_{12}\);
 \item una cifra de \(\alpha\) o una referencia metrológica interviene en
       cualquiera de las flechas de \eqref{p3-83-c25-s4:eq:composicion-upstream-u-k}.
\end{enumerate}
\end{remark}


% Unidad U017. Cierre afín, prolongación y semántica de frontera de alpha.
% Redacción autónoma desde la autoridad material 1063 y su sucesor V2.

\section{Prolongación coinductiva del cociclo dodecafásico y compatibilidad de la frontera \texorpdfstring{\(662/663\)}{662/663}}
\label{p3-83-c25-s5:chap:acarreo-alpha-frontera-662-663}

\subsection{Dominio aritmético del cierre}

La construcción recibe cuatro coordenadas dodecafásicas obtenidas en las
unidades precedentes:
\[
\begin{aligned}
 P={}&(141,592,653,589,793,238,462,643,383,279,502,884),\\
 E={}&(718,281,828,459,045,235,360,287,471,352,662,497),\\
 \Phi={}&(618,033,988,749,894,848,204,586,834,365,638,117),\\
 K={}&(234,543,140,729,659,824,621,058,914,794,146,601).
\end{aligned}
\]
Los tres primeros vectores son las doce primeras tríadas de las lecturas
arquimedianas fraccionarias de \(\pi\), \(e\) y \(\varphi\). El cuarto es
la coordenada dodecafásica del estado terminal enriquecido del Topological
Phase Kernel, construida antes de abrir la operación de acarreo.

Fijamos la base
\[
 B:=10^3=1000
\]
y la concatenación
\begin{equation}
 \iota(v_1,\ldots,v_{12})
 :=\sum_{m=1}^{12}v_mB^{12-m}.
 \label{p3-83-c25-s5:eq:concatenacion-base-mil-alpha}
\end{equation}
\Needspace{8\baselineskip}
En particular,
\[
\begin{aligned}
 \iota(P)&=141592653589793238462643383279502884,\\
 \iota(E)&=718281828459045235360287471352662497,\\
 \iota(\Phi)&=618033988749894848204586834365638117,\\
 \iota(K)&=234543140729659824621058914794146601.
\end{aligned}
\]

La operación estudiada en esta unidad es el mapa
\begin{equation}
 \mathfrak C_{12}:
 (P,E,\Phi,K,c_{13})
 \longmapsto(A,C),
 \label{p3-83-c25-s5:eq:mapa-cierre-dodecafasico-alpha}
\end{equation}
donde \(A=(A_1,\ldots,A_{12})\) es la palabra de salida y
\[
 C=(c_1,\ldots,c_{13})
\]
es el sistema completo de trece cocientes euclídeos: doce cocientes
internos y una condición de frontera derecha.

\subsection{División euclídea orientada de derecha a izquierda}

Para \(m=12,11,\ldots,1\), se define
\begin{equation}
\begin{aligned}
 s_m&:=P_m+E_m-\Phi_m-K_m+c_{m+1},\\
 A_m&:=s_m\bmod B,
       \qquad 0\leq A_m<B,\\
 c_m&:=\frac{s_m-A_m}{B}
      =\left\lfloor\frac{s_m}{B}\right\rfloor .
\end{aligned}
\label{p3-83-c25-s5:eq:recurrencia-corta-alpha}
\end{equation}
\Needspace{9\baselineskip}
La segunda igualdad para \(c_m\) es válida también cuando \(s_m<0\),
porque se utiliza división euclídea con resto en
\(\{0,\ldots,B-1\}\). Por ello,
\[
 -431=569-1000,
 \qquad
 -717=283-1000,
 \qquad
 -1200=800-2\cdot1000,
\]
y los préstamos negativos no introducen ambigüedad.

\begin{theorem}[Existencia y unicidad del acarreo orientado]
\label{p3-83-c25-s5:thm:unicidad-acarreo-orientado-alpha}
Fijados \(P,E,\Phi,K\) y el acarreo terminal \(c_{13}\), la recurrencia
\eqref{p3-83-c25-s5:eq:recurrencia-corta-alpha} determina de manera única
\[
 (A_1,\ldots,A_{12};c_1,\ldots,c_{12}).
\]
\end{theorem}

\begin{proof}
La fila \(m=12\) queda determinada por la división euclídea de
\[
 P_{12}+E_{12}-\Phi_{12}-K_{12}+c_{13}
\]
entre \(B\). Conocido \(c_{12}\), la fila \(m=11\) queda determinada por
la misma propiedad. La inducción descendente alcanza \(m=1\). En cada
paso, la unicidad del cociente y del resto euclídeos impide dos salidas
distintas.
\end{proof}

\subsection{Ventana dodecafásica finita}

La ventana finita se define mediante la condición terminal
\[
 c_{13}^{\rm fin}=0.
\]
La evaluación completa es la siguiente. La tabla se imprime en orden
creciente, aunque debe leerse causalmente desde la última fila hacia la
primera.

\begingroup
\scriptsize
\setlength{\tabcolsep}{3.2pt}
\begin{longtable}{@{}r|rrrr|r|r|r|r@{}}
\caption{Las doce divisiones euclídeas y los trece acarreos de la ventana
dodecafásica finita.}
\label{p3-83-c25-s5:tab:acarreo-alpha-finito-autonomo}\\
\toprule
\(m\)&\(P_m\)&\(E_m\)&\(\Phi_m\)&\(K_m\)&
\(c_{m+1}\)&\(s_m\)&\(A_m\)&\(c_m\)\\
\midrule
\endfirsthead
\toprule
\(m\)&\(P_m\)&\(E_m\)&\(\Phi_m\)&\(K_m\)&
\(c_{m+1}\)&\(s_m\)&\(A_m\)&\(c_m\)\\
\midrule
\endhead
1  &141&718&618&234& \(0\)& \(7\)&007& \(0\)\\
2  &592&281&033&543& \(0\)& \(297\)&297& \(0\)\\
3  &653&828&988&140& \(-1\)& \(352\)&352& \(0\)\\
4  &589&459&749&729& \(-1\)& \(-431\)&569& \(-1\)\\
5  &793&045&894&659& \(-2\)& \(-717\)&283& \(-1\)\\
6  &238&235&848&824& \(-1\)& \(-1200\)&800& \(-2\)\\
7  &462&360&204&621& \(0\)& \(-3\)&997& \(-1\)\\
8  &643&287&586&058& \(-1\)& \(285\)&285& \(0\)\\
9  &383&471&834&914& \(-1\)& \(-895\)&105& \(-1\)\\
10 &279&352&365&794& \(0\)& \(-528\)&472& \(-1\)\\
11 &502&662&638&146& \(0\)& \(380\)&380& \(0\)\\
12 &884&497&117&601& \(0\)& \(663\)&663& \(0\)\\
\bottomrule
\end{longtable}
\endgroup

La palabra completa de trece acarreos es
\[
 \boxed{
 C^{\rm fin}
 =(c_1,\ldots,c_{13})
 =(0,0,0,-1,-1,-2,-1,0,-1,-1,0,0,0).}
\]
La salida dodecafásica es
\[
 \boxed{
 A^{\rm fin}
 =(007,297,352,569,283,800,997,285,105,472,380,663).}
\]
Definimos el racional finito
\begin{equation}
 \boxed{
 \alpha_{12}^{\rm fin}
 :=10^{-36}\iota(A^{\rm fin})
 =\frac{7297352569283800997285105472380663}{10^{36}}.}
 \label{p3-83-c25-s5:eq:alpha-finita-autonoma}
\end{equation}
Por tanto,
\[
 \alpha_{12}^{\rm fin}
 =0.007297352569283800997285105472380663.
\]

\subsection{Reconocimiento metrológico posterior: CODATA 2018}

La comparación metrológica se efectúa sólo después de cerrar la recurrencia,
la frontera derecha y el racional anterior. La identidad entera da
\begin{equation}
 \boxed{
 \alpha_{12}^{\rm fin}
 =10^{-36}
 \left\lfloor\frac{10^{36}}{137.035999084}\right\rfloor.}
 \label{p3-83-c25-s5:eq:alpha-finita-codata-2018}
\end{equation}
En consecuencia,
\begin{equation}
 (\alpha_{12}^{\rm fin})^{-1}
 =137.03599908400000000000000000000001066588\ldots,
 \label{p3-83-c25-s5:eq:alpha-inversa-codata-2018}
\end{equation}
y reproduce exactamente la cadena central publicada por CODATA 2018
\cite{Tiesinga2021CODATA}:
\[
 \boxed{137.035999084.}
\]
La palabra \emph{exactamente} se refiere a la cadena central publicada, no
a una supuesta medición experimental de treinta y seis cifras. Las cifras
posteriores de \eqref{p3-83-c25-s5:eq:alpha-inversa-codata-2018} son consecuencias del
racional HMT ya fijado. CODATA no selecciona el estado, no determina (K),
no decide el sentido del acarreo y no interviene en la tabla dodecafásica.

\subsection{Identidad afín completa}

Cada fila de la tabla satisface exactamente
\begin{equation}
 P_m+E_m-\Phi_m-K_m+c_{m+1}
 =A_m+Bc_m.
 \label{p3-83-c25-s5:eq:identidad-local-acarreo-alpha}
\end{equation}
Introduciendo
\[
 C^-:=(c_1,\ldots,c_{12}),
 \qquad
 C^+:=(c_2,\ldots,c_{13}),
\]
las doce identidades locales se reúnen en
\begin{equation}
 \boxed{
 P+E-\Phi-K+C^+-BC^-=A.}
 \label{p3-83-c25-s5:eq:identidad-afin-vectorial-alpha}
\end{equation}
El término
\[
 \partial_BC:=C^+-BC^-
\]
es el cociclo de acarreo de la cadena orientada. En consecuencia,
\(P+E-\Phi-K=A\) no es una igualdad coordenada en
\(\mathbb Z^{12}\): sólo se vuelve correcta al incorporar
\(\partial_BC\).

\begin{theorem}[Telescopado de los trece acarreos]
\label{p3-83-c25-s5:thm:telescopado-acarreos-alpha}
Para cualquier condición terminal \(c_{13}\), la recurrencia satisface
\begin{equation}
 \boxed{
 \iota(P)+\iota(E)-\iota(\Phi)-\iota(K)-\iota(A)
 =B^{12}c_1-c_{13}.}
 \label{p3-83-c25-s5:eq:telescopado-general-alpha}
\end{equation}
En la ventana finita, \(c_1=c_{13}=0\), y por tanto
\begin{equation}
 \boxed{
 \iota(P)+\iota(E)-\iota(\Phi)-\iota(K)
 =\iota(A^{\rm fin}).}
 \label{p3-83-c25-s5:eq:identidad-entera-alpha-finita}
\end{equation}
\end{theorem}

\begin{proof}
Multiplicamos \eqref{p3-83-c25-s5:eq:identidad-local-acarreo-alpha} por
\(B^{12-m}\) y sumamos sobre \(m=1,\ldots,12\). La contribución de los
acarreos es
\[
 \sum_{m=1}^{12}
 (Bc_m-c_{m+1})B^{12-m}.
\]
Las componentes interiores se cancelan:
\begin{align*}
 \sum_{m=1}^{12}(Bc_m-c_{m+1})B^{12-m}
 &=\sum_{m=1}^{12}c_mB^{13-m}
   -\sum_{m=1}^{12}c_{m+1}B^{12-m}\\
 &=B^{12}c_1-c_{13}.
\end{align*}
Esto prueba \eqref{p3-83-c25-s5:eq:telescopado-general-alpha}. La especialización
\(c_1=c_{13}=0\) da \eqref{p3-83-c25-s5:eq:identidad-entera-alpha-finita}.
\end{proof}

\subsection{Inversión exacta del sello \texorpdfstring{\(K\)}{K}}

La recurrencia local es reversible. Despejando \(K_m\) en
\eqref{p3-83-c25-s5:eq:identidad-local-acarreo-alpha},
\begin{equation}
 \boxed{
 K_m=P_m+E_m-\Phi_m+c_{m+1}-A_m-Bc_m.}
 \label{p3-83-c25-s5:eq:inversion-local-k-desde-alpha}
\end{equation}
Por concatenación,
\begin{equation}
 \boxed{
 \iota(K)
 =\iota(P)+\iota(E)-\iota(\Phi)-\iota(A)
  -B^{12}c_1+c_{13}.}
 \label{p3-83-c25-s5:eq:inversion-concatenada-k-desde-alpha}
\end{equation}
En la frontera finita,
\[
 \iota(K)
 =\iota(P)+\iota(E)-\iota(\Phi)-\iota(A^{\rm fin}).
\]

La reversibilidad algebraica no invierte el orden causal. El orden efectivo
es
\[
 x_{\rm term}
 \longrightarrow K
 \longrightarrow(P,E,\Phi,K,c_{13})
 \longrightarrow(A,C).
\]
La fórmula \eqref{p3-83-c25-s5:eq:inversion-local-k-desde-alpha} es un control de
conmutatividad: demuestra que la salida y los acarreos recuperan el sello
que ya estaba fijado. No constituye una regla para seleccionar
retrospectivamente \(K\) a partir de un valor deseado de \(\alpha\).

\subsection{Prolongación coinductiva del sello dodecafásico y recurrencia inversa del acarreo en \(\mathcal C=\{-2,-1,0,1,2\}\)}

Sean
\[
 (P_k)_{k\geq1},
 \qquad
 (E_k)_{k\geq1},
 \qquad
 (\Phi_k)_{k\geq1}
\]
las secuencias infinitas de tríadas producidas por las tres ramas
coinductivas, y sea
\[
 \kappa(k):=1+((k-1)\bmod12).
\]
La prolongación del sello es periódica:
\[
 K_k^{\rm per}:=K_{\kappa(k)}.
\]
Para una profundidad finita \(N\) y una condición terminal
\(t=c_{N+1}\), la recurrencia es
\begin{equation}
\begin{aligned}
 s_k^{(N,t)}
 &=P_k+E_k-\Phi_k-K_{\kappa(k)}+c_{k+1}^{(N,t)},\\
 A_k^{(N,t)}
 &=s_k^{(N,t)}\bmod1000,\\
 c_k^{(N,t)}
 &=\left\lfloor\frac{s_k^{(N,t)}}{1000}\right\rfloor,
 \qquad k=N,\ldots,1.
\end{aligned}
\label{p3-83-c25-s5:eq:recurrencia-larga-alpha}
\end{equation}

El conjunto de acarreos terminales
\[
 \mathcal C:=\{-2,-1,0,1,2\}
\]
es invariante para las entradas certificadas: para cada tríada y cada
\(c_{k+1}\in\mathcal C\), el cociente \(c_k\) vuelve a pertenecer a
\(\mathcal C\). No se elige una condición terminal favorable. Se propagan
las cinco condiciones y sólo se publica el prefijo común a todas ellas.

La ejecución de profundidad certificada emplea
\[
\begin{array}{c|r}
\text{cantidad}&\text{valor}\\ \hline
\text{bloques ternarios por canal}&3501\\
\text{trits por canal}&3501\cdot6=21006\\
\text{cifras decimales de guarda}&10020\\
\text{tríadas de guarda}&3340\\
\text{condiciones terminales propagadas}&5\\
\text{tríadas comunes}&3339\\
\text{cifras comunes}&10017\\
\text{cifras publicadas}&10000
\end{array}
\]
Las cinco ramas terminales coalescen antes del prefijo publicado. En
particular, ninguna condición terminal extraída de un valor objetivo de
\(\alpha\) interviene en la salida.

\begin{theorem}[Compatibilidad de las salidas estabilizadas]
\label{p3-83-c25-s5:thm:compatibilidad-prolongacion-alpha}
Sean \(A^{[N]}\) y \(A^{[N']}\), con \(N'>N\), dos ejecuciones cuyas
cinco condiciones terminales han coalescido antes de la posición \(N\).
Entonces el truncamiento de \(A^{[N']}\) a sus primeros \(N\) bloques
coincide con \(A^{[N]}\). Los prefijos estabilizados definen una única
sección en el límite inverso de las palabras finitas.
\end{theorem}

\begin{proof}
La recurrencia es determinista de derecha a izquierda una vez conocido el
acarreo que entra en una posición. La coalescencia afirma que todas las
condiciones remotas producen el mismo acarreo antes de la frontera del
prefijo. Desde ese punto, las mismas tríadas \(P_k,E_k,\Phi_k\) y el mismo
sello periódico producen, por unicidad euclídea, los mismos restos y
cocientes en todas las posiciones anteriores. Esto es exactamente la
conmutatividad con el truncamiento.
\end{proof}

\subsection{Frontera procedente de la cola remota}

La prolongación determina en la frontera visible
\[
 c_{13}^{\infty}=-1.
\]
Los doce acarreos interiores siguen siendo los mismos que en la ventana
finita:
\[
 \boxed{
 C_{\leq13}^{\infty}
 =(0,0,0,-1,-1,-2,-1,0,-1,-1,0,0,-1).}
\]
Sólo cambia la información que entra desde la decimotercera tríada. Las
dos divisiones relevantes son
\[
\begin{aligned}
 s_{12}^{\infty}
 &=884+497-117-601-1=662,\\
 A_{12}^{\infty}&=662,
 \qquad c_{12}^{\infty}=0,
\end{aligned}
\]
y
\[
\begin{aligned}
 s_{13}^{\infty}
 &=197+757-720-234-1=-1,\\
 A_{13}^{\infty}&=999,
 \qquad c_{13}^{\infty}=-1.
\end{aligned}
\]
Por tanto, el comienzo de la palabra prolongada es
\[
 \boxed{
 A^{\infty}
 =(007,297,352,569,283,800,997,285,105,472,380,662,999,\ldots).}
\]
La correspondiente expansión comienza
\begin{equation}
 \boxed{
 \alpha_{\infty}
 =0.007297352569283800997285105472380662999564172539642\ldots.}
 \label{p3-83-c25-s5:eq:alpha-infinita-autonoma}
\end{equation}

La diferencia entre ambas fronteras es local y exacta:
\[
 c_{13}^{\rm fin}-c_{13}^{\infty}=1,
 \qquad
 A_{12}^{\rm fin}-A_{12}^{\infty}=1,
\]
mientras
\[
 A_m^{\rm fin}=A_m^\infty,
 \qquad 1\leq m\leq11.
\]

Aplicando el teorema de telescopado a los doce primeros bloques de la
prolongación,
\[
 \iota(P)+\iota(E)-\iota(\Phi)-\iota(K)
 -\iota(A^\infty_{\leq12})
 =-c_{13}^{\infty}=1.
\]
Luego
\begin{equation}
 \boxed{
 \iota(A^{\rm fin})
 =\iota(A^\infty_{\leq12})+1.}
 \label{p3-83-c25-s5:eq:relacion-unidad-fronteras-alpha}
\end{equation}
La terminación \(663\) no contradice la terminación \(662\): la primera
es el representante finito obtenido al clausurar la ventana con
\(c_{13}=0\); la segunda conserva el acarreo \(-1\) procedente de la cola
remota.

\begin{figure}[htbp]
\centering
\begin{tikzpicture}[
  >=Stealth,
  node distance=12mm and 16mm,
  box/.style={draw,rounded corners=1.5mm,align=center,
    inner xsep=3.2mm,inner ysep=2.5mm},
  arr/.style={->,thick}
]
\node[box] (shared) {datos compartidos\\
 \(P,E,\Phi,K\)};
\node[box,below left=of shared] (finite)
 {frontera finita\\\(c_{13}=0\)};
\node[box,below right=of shared] (remote)
 {cola remota\\\(c_{13}=-1\)};
\node[box,below=of finite] (a663)
 {\(A_{12}=663\)\\clausura finita};
\node[box,below=of remote] (a662)
 {\(A_{12}=662\), \(A_{13}=999\)\\prefijo prolongado};
\draw[arr] (shared)--(finite);
\draw[arr] (shared)--(remote);
\draw[arr] (finite)--(a663);
\draw[arr] (remote)--(a662);
\draw[<->,densely dashed] (a663)--node[below,align=center]
 {diferencia entera exacta \(1\)} (a662);
\end{tikzpicture}
\caption{Bifurcación de frontera de una misma recurrencia. No se comparan
dos generadores distintos: cambia únicamente el dato que entra desde la
derecha.}
\label{p3-83-c25-s5:fig:bifurcacion-frontera-alpha}
\end{figure}

\subsection{Truncamiento compatible y redondeo certificado a \(36\) cifras}

Para \(x\geq0\), definimos
\[
 \operatorname{Tr}_{36}(x)
 :=10^{-36}\left\lfloor10^{36}x\right\rfloor
\]
y
\[
 \operatorname{Rd}_{36}(x)
 :=10^{-36}\left\lfloor10^{36}x+\frac12\right\rfloor.
\]
Sea
\[
 a:=7297352569283800997285105472380662.
\]
La expansión \eqref{p3-83-c25-s5:eq:alpha-infinita-autonoma} puede escribirse como
\[
 \alpha_\infty=\frac{a+\rho}{10^{36}},
 \qquad
 \rho=0.999564172539642709\ldots,
\]
por lo que \(0.999<\rho<1\).

\begin{theorem}[Resolución exacta de la frontera \(662/663\)]
\label{p3-83-c25-s5:thm:frontera-alpha-662-663}
La prolongación coinductiva satisface
\[
 \boxed{
 \operatorname{Tr}_{36}(\alpha_\infty)
 =0.007297352569283800997285105472380662,}
\]
mientras su redondeo a treinta y seis cifras satisface
\[
 \boxed{
 \operatorname{Rd}_{36}(\alpha_\infty)
 =0.007297352569283800997285105472380663
 =\alpha_{12}^{\rm fin}.}
\]
\end{theorem}

\begin{proof}
Como \(0\leq\rho<1\),
\[
 \left\lfloor10^{36}\alpha_\infty\right\rfloor=a.
\]
Esto da el truncamiento terminado en \(662\). Como \(\rho>1/2\),
\[
 \left\lfloor10^{36}\alpha_\infty+\frac12\right\rfloor=a+1,
\]
que termina en \(663\). Por \eqref{p3-83-c25-s5:eq:alpha-finita-autonoma},
\((a+1)10^{-36}\) es exactamente \(\alpha_{12}^{\rm fin}\).
\end{proof}

Además,
\[
 0<\alpha_{12}^{\rm fin}-\alpha_\infty
 =\frac{1-\rho}{10^{36}}<10^{-39}.
\]
La diferencia no es un error aritmético de la tabla: cuantifica la
diferencia entre dos operadores de frontera distintos.

\subsection{Aislamiento causal respecto de \texorpdfstring{\(\alpha\)}{alpha} y de los datos metrológicos}

Para profundidad \(N\), el dominio completo de la operación es
\[
 \mathscr D_N
 =\bigl((P_k)_{k=1}^{N},(E_k)_{k=1}^{N},
        (\Phi_k)_{k=1}^{N},K,c_{N+1}\bigr).
\]
El mapa de acarreo
\[
 \mathfrak C_N:\mathscr D_N
 \longrightarrow
 \{0,\ldots,999\}^{N}\times\mathcal C^{N}
\]
está definido exclusivamente por \eqref{p3-83-c25-s5:eq:recurrencia-larga-alpha}. Su
dominio no contiene un valor objetivo de \(\alpha\), su inversa, una tabla
metrológica, una cadena decimal objetivo ni un acarreo terminal elegido por
comparación.

\begin{theorem}[No interferencia del blanco metrológico]
\label{p3-83-c25-s5:thm:no-interferencia-metrologica-alpha}
Sea \(\mathscr Y\) cualquier espacio de datos externos que contenga valores
metrológicos o tablas de comparación. La extensión formal
\[
 \widetilde{\mathfrak C}_N:
 \mathscr D_N\times\mathscr Y
 \longrightarrow
 \{0,\ldots,999\}^{N}\times\mathcal C^{N},
 \qquad
 \widetilde{\mathfrak C}_N(d,y):=\mathfrak C_N(d),
\]
factoriza por la proyección
\[
 \operatorname{pr}_{\mathscr D_N}:
 \mathscr D_N\times\mathscr Y\longrightarrow\mathscr D_N.
\]
Por tanto, para cualesquiera \(y,y'\in\mathscr Y\),
\[
 \widetilde{\mathfrak C}_N(d,y)
 =\widetilde{\mathfrak C}_N(d,y').
\]
Ninguna intervención sobre un valor de \(\alpha\) o sobre una tabla
metrológica altera \(K\), los cocientes, los restos ni el prefijo producido.
\end{theorem}

\begin{proof}
La afirmación resulta de
\[
 \widetilde{\mathfrak C}_N
 =\mathfrak C_N\circ\operatorname{pr}_{\mathscr D_N}.
\]
Todas las operaciones de \(\mathfrak C_N\) son sumas, restas y divisiones
euclídeas de los elementos de \(\mathscr D_N\). Ninguna coordenada de
\(\mathscr Y\) aparece en ellas.
\end{proof}

El orden causal completo es
\[
 \text{ramas enriquecidas de }\pi,e,\varphi
 \longrightarrow(P,E,\Phi),
\]
\[
 x_{\rm term}\longrightarrow K,
\]
\[
 (P,E,\Phi,K,\text{frontera})
 \longrightarrow(A,C)
 \longrightarrow\alpha_\infty,
\]
\[
 \alpha_\infty
 \longrightarrow\text{comparación externa}.
\]
La última flecha no puede invertirse para seleccionar ninguna de las
anteriores.

\subsection{Obstrucciones exactas de la prolongación: frontera, truncamiento y aislamiento causal}

\begin{remark}[Certificado y control de validez: Refutación del cierre afín y de su prolongación]
El bloque queda refutado en su dominio si ocurre cualquiera de los hechos
siguientes:
\begin{enumerate}
 \item alguna fila de la \cref{p3-83-c25-s5:tab:acarreo-alpha-finito-autonomo}
       incumple \(s_m=A_m+1000c_m\);
 \item el vector publicado de trece acarreos no reproduce las doce filas;
 \item falla la identidad telescópica
       \[
       \iota(P)+\iota(E)-\iota(\Phi)-\iota(K)-\iota(A)
       =1000^{12}c_1-c_{13};
       \]
 \item la inversión \eqref{p3-83-c25-s5:eq:inversion-local-k-desde-alpha} no recupera
       las doce coordenadas de \(K\);
 \item la propagación remota produce \(c_{13}\neq-1\) para el prefijo
       certificado;
 \item las cinco condiciones terminales de
       \(\mathcal C=\{-2,-1,0,1,2\}\) no coalescen antes de las cifras
       publicadas;
 \item la tríada siguiente a \(662\) no es \(999\), en cuyo caso deja de
       estar demostrada la identificación del redondeo con \(663\);
 \item el generador abre una expansión objetivo de \(\alpha\), su inversa,
       una tabla metrológica o una condición terminal elegida para forzar
       la coincidencia.
\end{enumerate}
\end{remark}

Quedan así demostradas dos identidades de frontera distintas:
\[
 \boxed{
 663=\text{clausura finita con }c_{13}=0
     =\text{redondeo a 36 cifras de }\alpha_\infty,}
\]
\[
 \boxed{
 662=\text{truncamiento a 36 cifras de }\alpha_\infty
 \quad\text{con}\quad c_{13}=-1,\ A_{13}=999.}
\]
El resultado metrológico central queda, finalmente, publicado en la misma
residencia que la derivación:
\begin{equation}
 \boxed{
 \begin{gathered}
  (\alpha_{12}^{\rm fin})^{-1}
  =137.03599908400000000000000000000001066588\ldots,\\
  \text{cadena central CODATA 2018:}\qquad137.035999084.
 \end{gathered}}
 \label{p3-83-c25-s5:eq:resultado-final-alpha-codata-2018}
\end{equation}
La implicación expresa comparación posterior con la cadena central
publicada; no introduce esa cadena en el constructor.
